% Part: lambda-calculus
% Chapter: church-rosser
% Section: definitions-and-properties

\documentclass[../../../include/open-logic-section]{subfiles}

\begin{document}

\olfileid{lam}{cr}{dap}

\olsection{定义与性质}

本章我们介绍 Church--Rosser（丘奇—罗瑟）性质的概念及其一些常见性质。

\begin{defn}[Church--Rosser 性质, CR]
  项上的关系 $\xredone$ 满足\emph{Church--Rosser 性质}，当且仅当：只要 $M \xredone P$ 且 $M \xredone Q$，就存在某个项~$N$，使得 $P \xredone N$ 且 $Q \xredone N$。
\end{defn}

我们可以把λ演算看作一个计算模型，其中处于范式的项就是“值”，而项上的可归约关系就是“计算规则”。Church--Rosser 性质表明：当计算存在不止一种进行方式时，表达式的值仍然是唯一的。

以初等代数为例，计算 $4 \times (1+2) + 3$ 的方式不止一种。它可以先归约为 $4 \times 3+3$（如果先把 $1+2$ 算成~$3$），也可以先归约为 $4 \times 1+4 \times 2+3$（如果先用分配律算 $4 \times (1+2)$）。然而，这两者都可以进一步归约为 $12+3$。

如果我们令 $\xredone$ 为 β-归约，很容易看出 Church--Rosser 性质的一个推论是：如果一个项有范式，那么该范式是唯一的。这是因为，假设 $M$ 可以归约为 $P$ 和 $Q$，且二者都是范式。根据 Church--Rosser 性质，存在某个 $N$ 使得 $P$ 和 $Q$ 都能归约到该 $N$。由假设，$P$ 和 $Q$ 都是范式，它们归约到 $N$ 的过程只能是平凡归约，即 $P$、$Q$ 和 $N$ 是同一项。这就为我们谈论一个项的\emph{那个}范式提供了依据。

因此，在将λ演算视为计算模型时，项的范式可以被看作从该项开始计算所得的“最终结果”。上述推论意味着，如果计算存在最终结果，那么该结果是唯一的，正如计算 $4 \times (1+2)+3$ 的结果只有一个，即~$15$。

\begin{thm}
\ollabel{thm:str}
  如果一个关系 $\xredone$ 满足 Church--Rosser 性质，并且 $\xred$ 是包含 $\xredone$ 的最小传递关系，那么 $\xred$ 也满足 Church--Rosser 性质。
\end{thm}

\begin{proof}
  假设
  \begin{align*}
    M & \xredone P_1 \xredone \dots \xredone P_m \text{ 且}\\
    M & \xredone Q_1 \xredone \dots \xredone Q_n.
  \end{align*}。
  我们将通过构造一个高度为 $m + 1$、宽度为 $n + 1$ 的项网格 $N$ 来证明本定理。我们用 $N_{i,j}$ 表示第 $i$ 行第 $j$ 列的项。

  我们这样构造 $N$，使得 $N_{i,j} \xredone N_{i+1,j}$ 且
  $N_{i,j} \xredone N_{i,j+1}$。其定义如下：
  \begin{align*}
    N_{0,0} &= M \\
    N_{i,0} &= P_i && \text{若 } 1 \le i \le m \\
    N_{0,j} &= Q_j && \text{若 } 1 \le j \le n \\
  \intertext{否则：}
    N_{i,j} &= R
  \end{align*}
  其中 $R$ 是一个满足 $N_{i-1,j} \xredone R$ 且 $N_{i,j-1}
  \xredone R$ 的项。根据 $\xredone$ 的 Church--Rosser 性质，这样的项总是存在的。

  现在我们有 $N_{m,0} \xredone \dots \xredone N_{m,n}$ 和 $N_{0,n}
  \xredone \dots \xredone N_{m,n}$。注意 $N_{m,0}$ 就是 $P$，而 $N_{0,n}$ 就是 $Q$。根据 $\xred$ 的定义，定理得证。
\end{proof}

\end{document}